% answers.tex -- marking criteria and sample answers for questions.tex (Module 5 short answer).
% NESA marking-guideline style: a criteria table for each part, then a sample answer.
% Every number: checks/m5-short-answer/numbers.py
% Build: TEXINPUTS=.: xelatex answers.tex (twice)
\documentclass[12pt]{article}
\usepackage{quiz}
\hypersetup{pdftitle={Short Answer Questions - Module 5 - Advanced Mechanics - Marking guidelines}}
% links between the two PDF files are not possible, so question headings are plain text here
\renewcommand\hyperlink[2]{#2}
\newcommand\anstopic[1]{\par\Needspace{8\baselineskip}\vspace{4mm}%
  \begin{center}{\headfont\bfseries\fontsize{14}{18}\selectfont #1\par}\end{center}\vspace{2pt}}

\begin{document}
\quiztitle{Year 12 Physics}{Short Answer Questions - Module 5 - Advanced Mechanics:\\
  Marking Guidelines and Sample Answers}

Data used: $g = \qty{9.8}{m.s^{-2}}$, $G = \qty{6.67e-11}{N.m^2.kg^{-2}}$,
$M_\mathrm{E} = \qty{6.0e24}{kg}$, $r_\mathrm{E} = \qty{6.371e6}{m}$. In calculations, a correct
answer from a correct method with an earlier error carried forward still earns the later marks.

% =====================================================================
\anstopic{Topic 1 - Projectile Motion}

% ---- Q1
\ansq
\anspart{a}
\begin{crit}
  \critrow{Correctly calculates the maximum height}{2}
  \critrow{Calculates the vertical component of the initial velocity OR uses a correct equation for
    the maximum height}{1}
\end{crit}
\sample
$u_y = 28\sin 35^\circ = \qty{16.1}{m.s^{-1}}$. At the maximum height $v_y = 0$:
\[ v_y^2 = u_y^2 + 2as \quad\Rightarrow\quad 0 = 16.06^2 - 2(9.8)s \quad\Rightarrow\quad s = \qty{13.2}{m} \]
\anspart{b}
\begin{crit}
  \critrow{Correctly calculates the range}{2}
  \critrow{Calculates the time of flight OR the horizontal component of the velocity}{1}
\end{crit}
\sample
Time of flight $t = \dfrac{2u_y}{g} = \dfrac{2 \times 16.06}{9.8} = \qty{3.28}{s}$.
$u_x = 28\cos 35^\circ = \qty{22.94}{m.s^{-1}}$. Range $= u_x t = 22.94 \times 3.278 = \qty{75.2}{m}$.

% ---- Q2
\ansq
\anspart{a}
\begin{crit}
  \critrow{Correctly calculates the time}{2}
  \critrow{Uses a correct equation with an initial vertical velocity of zero}{1}
\end{crit}
\sample
Vertically $u_y = 0$: $s = \frac{1}{2}gt^2 \Rightarrow 30 = \frac{1}{2}(9.8)t^2 \Rightarrow
t = \sqrt{\dfrac{60}{9.8}} = \qty{2.47}{s}$.
\anspart{b}
\begin{crit}
  \critrow{Correctly calculates the magnitude AND direction of the final velocity}{3}
  \critrow{Correctly calculates the magnitude OR calculates $v_y$ and the direction}{2}
  \critrow{Calculates the final vertical velocity}{1}
\end{crit}
\sample
$v_y = gt = 9.8 \times 2.47 = \qty{24.2}{m.s^{-1}}$ downwards, and $v_x = \qty{15}{m.s^{-1}}$.
\[ v = \sqrt{15^2 + 24.2^2} = \qty{28.5}{m.s^{-1}} \qquad
   \theta = \tan^{-1}\!\left(\frac{24.2}{15}\right) = 58^\circ \text{ below the horizontal} \]

% ---- Q3
\ansq
\begin{crit}
  \critrow{Explains the horizontal motion (no force, constant velocity) AND the vertical motion
    (constant force, constant acceleration) using Newton's laws, AND links them to the parabolic
    shape}{3}
  \critrow{Explains the horizontal and vertical motions using Newton's laws}{2}
  \critrow{Provides some relevant information}{1}
\end{crit}
\sample
Ignoring air resistance, the only force on the projectile is its weight, which is vertical and
constant.
\begin{itemize}
  \item Horizontally there is no force, so by Newton's first law the horizontal velocity is constant:
        $x = u_x t$.
  \item Vertically there is a constant force, so by Newton's second law there is a constant
        acceleration $g$ downwards: $y = u_y t - \frac{1}{2}gt^2$.
\end{itemize}
Substituting $t = \dfrac{x}{u_x}$ gives $y = \dfrac{u_y}{u_x}x - \dfrac{g}{2u_x^2}x^2$. This is a
quadratic in $x$, so the path is a parabola.

% ---- Q4
\ansq
\begin{crit}
  \critrow{$v_x$: horizontal straight line at a constant positive value ($u\cos\theta$, or $u_x$) from 0
    to $T$}{1}
  \critrow{$v_y$: straight line with a constant negative gradient}{1}
  \critrow{$v_y$ starts at $+u\sin\theta$ and ends at $-u\sin\theta$ (or $+u_y$ and $-u_y$)}{1}
  \critrow{$v_y$ crosses zero at $\frac{T}{2}$ OR the gradient is labelled $-g$ ($\qty{-9.8}{m.s^{-2}}$)}{1}
\end{crit}
\sample
\begin{center}
  \begin{tikzpicture}
    \draw[axisline] (0,0) -- (5.2,0) node[right] {$t$};
    \draw[axisline] (0,-2.2) -- (0,2.4) node[above] {$v_x$};
    \node[left, font=\small] at (0,0) {0};
    \draw (4.4,0.08) -- (4.4,-0.08) node[below, font=\small] {$T$};
    \draw[line width=1.2pt] (0,1.2) -- (4.4,1.2) node[above left, font=\small] {constant $= u\cos\theta$};
    \node[left, font=\small] at (0,1.2) {$u\cos\theta$};
  \end{tikzpicture}\hspace{14mm}%
  \begin{tikzpicture}
    \draw[axisline] (0,0) -- (5.2,0) node[right] {$t$};
    \draw[axisline] (0,-2.2) -- (0,2.4) node[above] {$v_y$};
    \node[left, font=\small] at (0,0) {0};
    \draw (4.4,0.08) -- (4.4,-0.08) node[below, font=\small] {$T$};
    \draw[line width=1.2pt] (0,1.7) -- (4.4,-1.7);
    \node[left, font=\small] at (0,1.7) {$u\sin\theta$}; \node[left, font=\small] at (0,-1.7) {$-u\sin\theta$};
    \draw[guide] (0,-1.7) -- (4.4,-1.7);
    \node[above right, font=\small] at (2.2,0) {$\frac{T}{2}$};
    \node[right, font=\small] at (2.6,0.9) {gradient $= -g$};
  \end{tikzpicture}
\end{center}

% ---- Q5
\ansq
\anspart{a}
\begin{crit}
  \critrow{Plots all five points correctly AND draws an appropriate straight line of best fit}{3}
  \critrow{Plots at least four points correctly AND draws a line of best fit}{2}
  \critrow{Plots some points correctly}{1}
\end{crit}
\sample
\begin{center}
  \begin{tikzpicture}
    \begin{axis}[width=110mm, height=75mm, axis lines=left, xmin=0, xmax=3.5, ymin=0, ymax=1.6,
        xtick={0,0.5,...,3.5}, ytick={0,0.2,...,1.6}, grid=major, major grid style={black!25},
        tick label style={font=\small}, label style={font=\small},
        xticklabel style={/pgf/number format/fixed, /pgf/number format/precision=1},
        yticklabel style={/pgf/number format/fixed, /pgf/number format/precision=1},
        xlabel={Launch speed $u$ (\si{m.s^{-1}})}, ylabel={Range $R$ (\si{m})},
        scaled ticks=false, enlargelimits=false]
      \addplot[only marks, mark=x, mark size=3pt, thick] coordinates {(1.0,0.43) (1.5,0.63) (2.0,0.87) (2.5,1.06) (3.0,1.29)};
      \addplot[line width=0.8pt, domain=0:3.4] {0.43*x};
    \end{axis}
  \end{tikzpicture}
\end{center}
\anspart{b}
\begin{crit}
  \critrow{Calculates the gradient, relates it to the time of flight AND calculates a value for $g$
    consistent with the gradient}{3}
  \critrow{Calculates the gradient AND relates it to the time of flight (or to $g$)}{2}
  \critrow{Calculates the gradient OR states a correct relationship between $R$, $u$ and $g$}{1}
\end{crit}
\sample
Gradient $= \dfrac{1.29 - 0}{3.0 - 0} \approx \qty{0.43}{s}$ (accept 0.42 to 0.44).
\[ R = ut \text{ and } h = \tfrac{1}{2}gt^2 \;\Rightarrow\; R = u\sqrt{\frac{2h}{g}}, \text{ so the
   gradient} = t = \sqrt{\frac{2h}{g}} \]
\[ g = \frac{2h}{(\text{gradient})^2} = \frac{2 \times 0.90}{0.43^2} = \qty{9.7}{m.s^{-2}} \]
(Any value from about 9.3 to \qty{10.2}{m.s^{-2}} that matches the student's own gradient is
accepted.)
\anspart{c}
\begin{crit}
  \critrow{Explains a valid improvement AND how it makes the results more accurate}{2}
  \critrow{Identifies a valid improvement}{1}
\end{crit}
\sample
Put carbon paper on the floor where the ball lands, so the landing point is marked exactly. The range
can then be measured precisely from the mark, instead of estimating where the ball hit. This reduces
the error in each value of $R$, so the results are closer to the true values. (Other valid answers:
record the landing with a video camera; place the light gate right at the edge of the bench so $u$ is
the true launch speed; use a dense, small ball to reduce air resistance.)

% ---- Q6
\ansq
\begin{crit}
  \critrow{Correctly calculates the horizontal distance}{4}
  \critrow{Correctly calculates the time of flight}{3}
  \critrow{Sets up a correct vertical equation with correct signs}{2}
  \critrow{Resolves the initial velocity into components}{1}
\end{crit}
\sample
$u_x = 12\cos 40^\circ = \qty{9.19}{m.s^{-1}}$ and $u_y = 12\sin 40^\circ = \qty{7.71}{m.s^{-1}}$.
Take up as positive, so $s_y = \qty{-25}{m}$ and $a = \qty{-9.8}{m.s^{-2}}$:
\[ -25 = 7.71t - 4.9t^2 \;\Rightarrow\; 4.9t^2 - 7.71t - 25 = 0 \;\Rightarrow\;
   t = \frac{7.71 + \sqrt{7.71^2 + 4(4.9)(25)}}{2(4.9)} = \qty{3.18}{s} \]
Horizontal distance $= u_x t = 9.19 \times 3.18 = \qty{29.2}{m}$.

% =====================================================================
\anstopic{Topic 2 - Circular Motion}

% ---- Q7
\ansq
\anspart{a}
\begin{crit}
  \critrow{Identifies friction between the tyres and the road (acting towards the centre)}{1}
\end{crit}
\sample
The sideways friction force between the tyres and the road, acting towards the centre of the bend.
\anspart{b}
\begin{crit}
  \critrow{Relates a higher speed to a larger required centripetal force AND links this to the
    maximum friction force}{2}
  \critrow{Provides some relevant information}{1}
\end{crit}
\sample
The centripetal force needed is $F = \dfrac{mv^2}{r}$, so it increases with $v^2$. Friction has a
maximum value. At a higher speed the force needed can be larger than the maximum friction, so the
car cannot follow the curve: it skids outwards, along a straighter path.

% ---- Q8
\ansq
\anspart{a}
\begin{crit}
  \critrow{Correctly draws and labels the weight AND the normal force, AND shows the net force
    horizontally towards the centre}{2}
  \critrow{Correctly draws the two forces OR shows the correct direction of the net force}{1}
\end{crit}
\sample
\begin{center}
  \begin{tikzpicture}[scale=1.3]
    \draw[line width=0.6pt] (-0.5,0) -- (5.6,0);
    \draw[line width=0.8pt] (0,0) -- (5.2,{5.2*tan(18)});
    \draw (0.9,0) arc (0:18:0.9); \node[font=\small] at (1.15,0.17) {$\theta$};
    \begin{scope}[shift={({3.3*cos(18)},{3.3*sin(18)})}, rotate=18]
      \draw[fill=black!8, rounded corners=2pt] (-0.7,0.12) rectangle (0.7,0.8);
    \end{scope}
    \coordinate (C) at ({3.3*cos(18)-0.46*sin(18)},{3.3*sin(18)+0.46*cos(18)});
    \draw[force] (C) -- ++({-1.5*sin(18)},{1.5*cos(18)}) node[above, font=\small] {Normal force $N$};
    \draw[force] (C) -- ++(0,-1.4) node[below, font=\small] {Weight $mg$};
    \draw[->, line width=1.4pt, dash pattern=on 4pt off 2pt] (C) -- ++({-1.4*tan(18)},0) node[left, font=\small, align=center] {Net force\\(to the centre)};
  \end{tikzpicture}
\end{center}
\anspart{b}
\begin{crit}
  \critrow{Writes correct vertical and horizontal equations AND combines them to obtain the
    relationship}{3}
  \critrow{Writes correct vertical AND horizontal equations}{2}
  \critrow{Writes one correct equation OR resolves the normal force}{1}
\end{crit}
\sample
The car does not accelerate vertically, so the vertical forces balance: $N\cos\theta = mg$.
The horizontal component of $N$ is the centripetal force: $N\sin\theta = \dfrac{mv^2}{r}$.
Dividing the second equation by the first:
\[ \frac{N\sin\theta}{N\cos\theta} = \frac{mv^2/r}{mg} \quad\Rightarrow\quad \tan\theta = \frac{v^2}{rg} \]
\anspart{c}
\begin{crit}
  \critrow{Correctly calculates the design speed}{2}
  \critrow{Correct substitution into the relationship}{1}
\end{crit}
\sample
$v = \sqrt{rg\tan\theta} = \sqrt{80 \times 9.8 \times \tan 15^\circ} = \qty{14.5}{m.s^{-1}}$
(about \qty{52}{km.h^{-1}}).

% ---- Q9
\ansq
\anspart{a}
\begin{crit}
  \critrow{Correctly calculates the tension}{2}
  \critrow{Uses $T\cos\theta = mg$}{1}
\end{crit}
\sample
Vertically: $T\cos 30^\circ = mg$, so $T = \dfrac{0.25 \times 9.8}{\cos 30^\circ} = \qty{2.83}{N}$.
\anspart{b}
\begin{crit}
  \critrow{Correctly calculates the speed}{2}
  \critrow{Calculates the radius OR uses $T\sin\theta = \frac{mv^2}{r}$}{1}
\end{crit}
\sample
$r = 0.80\sin 30^\circ = \qty{0.40}{m}$. Horizontally: $T\sin 30^\circ = \dfrac{mv^2}{r}$, so
\[ v = \sqrt{\frac{rT\sin 30^\circ}{m}} = \sqrt{\frac{0.40 \times 2.83 \times 0.5}{0.25}} = \qty{1.50}{m.s^{-1}} \]
(or $v = \sqrt{rg\tan 30^\circ}$, which gives the same answer).

% ---- Q10
\ansq
\anspart{a}
\begin{crit}
  \critrow{Correctly calculates the minimum speed}{2}
  \critrow{States that at the minimum speed the tension is zero, so $mg = \frac{mv^2}{r}$}{1}
\end{crit}
\sample
At the minimum speed the tension at the top is zero, so the weight alone gives the centripetal
force: $mg = \dfrac{mv^2}{r}$, so $v = \sqrt{gr} = \sqrt{9.8 \times 0.60} = \qty{2.4}{m.s^{-1}}$.
\anspart{b}
\begin{crit}
  \critrow{Correctly calculates the tension}{2}
  \critrow{Uses $T - mg = \frac{mv^2}{r}$ (or calculates $\frac{mv^2}{r}$)}{1}
\end{crit}
\sample
At the bottom, tension acts up (to the centre) and weight acts down:
$T - mg = \dfrac{mv^2}{r}$, so
$T = 0.20 \times 9.8 + \dfrac{0.20 \times 6.0^2}{0.60} = 1.96 + 12.0 = \qty{14.0}{N}$.

% ---- Q11
\ansq
\anspart{a}
\begin{crit}
  \critrow{Correctly calculates the torque}{2}
  \critrow{Uses $\tau = rF\sin\theta$}{1}
\end{crit}
\sample
$\tau = rF\sin\theta = 0.35 \times 150 \times \sin 70^\circ = \qty{49}{N.m}$.
\anspart{b}
\begin{crit}
  \critrow{Correctly calculates the force (applied at \ang{90})}{1}
\end{crit}
\sample
The force is smallest when it is at \ang{90} to the handle: $F = \dfrac{\tau}{r} = \dfrac{65}{0.35}
= \qty{186}{N}$.
\anspart{c}
\begin{crit}
  \critrow{Relates torque to the distance from the pivot AND explains that the same force now gives
    a larger torque}{2}
  \critrow{States that a longer lever arm increases the torque}{1}
\end{crit}
\sample
Torque $\tau = rF\sin\theta$. The pipe increases $r$, the distance from the nut to where the force is
applied. So the same force gives a larger torque, and a smaller force is enough to reach the
\qty{65}{N.m} needed. For example, \qty{150}{N} at \ang{90} gives \qty{65}{N.m} when
$r = \frac{65}{150} = \qty{0.43}{m}$.

% ---- Q12
\ansq
\anspart{a}
\begin{crit}
  \critrow{Uses $F = \frac{mv^2}{r}$ to show that $F \propto v^2$ (with $m$ and $r$ constant) AND
    explains that this gives a straight line through the origin}{2}
  \critrow{Provides some relevant information}{1}
\end{crit}
\sample
$F = \dfrac{mv^2}{r}$. With $m$ and $r$ constant, $F \propto v^2$. So a graph of $F$ against $v^2$ is
a straight line through the origin, and its gradient can be used to find a quantity. A graph of $F$
against $v$ would be a curve (a parabola), which is hard to analyse.
\anspart{b}
\begin{crit}
  \critrow{Calculates the gradient, relates it to $\frac{m}{r}$ AND correctly calculates the mass}{3}
  \critrow{Calculates the gradient AND relates it to $\frac{m}{r}$}{2}
  \critrow{Calculates the gradient OR states that the gradient equals $\frac{m}{r}$}{1}
\end{crit}
\sample
Gradient $= \dfrac{3.0 - 0}{60 - 0} = 0.050$ (units: $\si{N}$ per $\si{m^2.s^{-2}}$, which is
\si{kg.m^{-1}}). Since $F = \dfrac{m}{r}v^2$, the gradient is $\dfrac{m}{r}$:
\[ m = \text{gradient} \times r = 0.050 \times 0.60 = \qty{0.030}{kg} \;(\qty{30}{g}) \]
\anspart{c}
\begin{crit}
  \critrow{Explains that the timing error (such as reaction time) is spread over 10 revolutions, so
    the percentage error in the period, and hence in $v$, is smaller}{2}
  \critrow{Provides some relevant information}{1}
\end{crit}
\sample
Starting and stopping the stopwatch has an error of about \qty{0.2}{s} (reaction time). This error is
about the same whether 1 or 10 revolutions are timed. Dividing the time for 10 revolutions by 10
makes the error in one period 10 times smaller, so the percentage error in $T$, and in
$v = \frac{2\pi r}{T}$, is much smaller.

% ---- Q13
\ansq
\begin{crit}
  \critrow{Explains that the net force is perpendicular to the velocity, so it does no work (or has
    no component along the velocity), so the speed and kinetic energy do not change}{3}
  \critrow{Gives two of these points}{2}
  \critrow{Gives one relevant point}{1}
\end{crit}
\sample
In uniform circular motion the net (centripetal) force always points to the centre, at \ang{90} to the
velocity. Work done is $W = Fs\cos\theta$, and with $\theta = 90^\circ$, $W = 0$. Since no work is
done on the object, its kinetic energy does not change (work--energy theorem). The force only changes
the \emph{direction} of the velocity, not the speed.

% =====================================================================
\anstopic{Topic 3 - Motion in Gravitational Fields}

% ---- Q14
\ansq
\anspart{a}
\begin{crit}
  \critrow{Correctly calculates the gravitational field strength}{2}
  \critrow{Correct substitution into $g = \frac{GM}{r^2}$}{1}
\end{crit}
\sample
$g = \dfrac{GM}{r^2} = \dfrac{(\num{6.67e-11})(\num{6.42e23})}{(\num{3.39e6})^2} = \qty{3.73}{N.kg^{-1}}$.
\anspart{b}
\begin{crit}
  \critrow{Correctly calculates the weight}{1}
\end{crit}
\sample
$W = mg = 70 \times 3.73 = \qty{261}{N}$.

% ---- Q15
\ansq
\anspart{a}
\begin{crit}
  \critrow{States that gravity provides the centripetal force, substitutes $v = \frac{2\pi r}{T}$
    AND rearranges correctly}{3}
  \critrow{Equates gravitational and centripetal force AND substitutes for $v$}{2}
  \critrow{Gives one correct step}{1}
\end{crit}
\sample
Gravity provides the centripetal force: $\dfrac{GMm}{r^2} = \dfrac{mv^2}{r}$. For one orbit,
$v = \dfrac{2\pi r}{T}$, so
\[ \frac{GMm}{r^2} = \frac{m}{r}\cdot\frac{4\pi^2 r^2}{T^2} = \frac{4\pi^2 mr}{T^2}
   \quad\Rightarrow\quad GM\,T^2 = 4\pi^2 r^3 \quad\Rightarrow\quad \frac{r^3}{T^2} = \frac{GM}{4\pi^2} \]
\anspart{b}
\begin{crit}
  \critrow{Correctly calculates the orbital radius}{2}
  \critrow{Converts the period to seconds AND substitutes correctly}{1}
\end{crit}
\sample
$T = 24.0 \times 3600 = \qty{86400}{s}$.
\[ r = \sqrt[3]{\frac{GMT^2}{4\pi^2}} = \sqrt[3]{\frac{(\num{6.67e-11})(\num{6.0e24})(86\,400)^2}{4\pi^2}}
   = \qty{4.23e7}{m} \]

% ---- Q16
\ansq
\begin{crit}
  \critrow{Compares altitude, speed AND period, with correct relationships between them, AND gives
    a correct use of each type of satellite}{4}
  \critrow{Compares most of these features correctly AND gives a use}{3}
  \critrow{Describes some features of each orbit}{2}
  \critrow{Gives one relevant feature or use}{1}
\end{crit}
\sample
\begin{itemize}
  \item \textbf{Altitude:} LEO satellites orbit about 200 to \qty{2000}{km} above the surface. GEO
        satellites orbit about \qty{36000}{km} above the equator.
  \item \textbf{Speed:} $v = \sqrt{\frac{GM}{r}}$, so the lower LEO satellites move faster (about
        \qty{7.8}{km.s^{-1}}) than GEO satellites (about \qty{3.1}{km.s^{-1}}).
  \item \textbf{Period:} $T^2 \propto r^3$, so LEO periods are short (about 90 minutes), while the
        GEO period is 24 hours, the same as Earth's rotation, in the same direction. A GEO satellite
        stays above the same point on the equator; a LEO satellite moves quickly across the sky.
  \item \textbf{Uses:} LEO satellites are close to Earth, so they are used for detailed imaging and
        remote sensing, and for the ISS. GEO satellites stay above one place, so they are used for
        communications (fixed dishes) and for weather monitoring of one region.
\end{itemize}
Similarity: in both orbits gravity provides the centripetal force, and the speed and period depend
only on the orbital radius.

% ---- Q17
\ansq
\anspart{a}
\begin{crit}
  \critrow{Correctly calculates the change in gravitational potential energy (including the
    sign or stating it is an increase)}{2}
  \critrow{Uses $U = -\frac{GMm}{r}$ for both orbits}{1}
\end{crit}
\sample
\[ \Delta U = \left(-\frac{GMm}{r_2}\right) - \left(-\frac{GMm}{r_1}\right)
   = GMm\left(\frac{1}{r_1} - \frac{1}{r_2}\right) \]
\[ \Delta U = (\num{6.67e-11})(\num{6.0e24})(500)\left(\frac{1}{\num{7.0e6}} - \frac{1}{\num{1.4e7}}\right) \]
$\Delta U = \qty[retain-explicit-plus]{+1.43e10}{J}$ (an increase).
\anspart{b}
\begin{crit}
  \critrow{Correctly calculates the change in kinetic energy (including the sign or stating it is a
    decrease)}{2}
  \critrow{Uses $K = \frac{GMm}{2r}$ (or calculates the orbital speeds)}{1}
\end{crit}
\sample
In a circular orbit $K = \dfrac{GMm}{2r}$, so
\[ \Delta K = \frac{GMm}{2}\left(\frac{1}{r_2} - \frac{1}{r_1}\right) = \qty{-7.15e9}{J} \text{ (a decrease)} \]
\anspart{c}
\begin{crit}
  \critrow{Correctly determines the minimum energy}{1}
\end{crit}
\sample
Energy supplied $= \Delta U + \Delta K = \num{1.43e10} - \num{7.15e9} = \qty{7.15e9}{J}$.

% ---- Q18
\ansq
\anspart{a}
\begin{crit}
  \critrow{States the escape condition (reaches infinity with zero kinetic energy, so the total
    energy is zero), applies conservation of energy AND derives the expression}{3}
  \critrow{Writes a correct energy-conservation equation}{2}
  \critrow{Gives one relevant point (for example, $U = -\frac{GMm}{r}$)}{1}
\end{crit}
\sample
To just escape, the object reaches a very large distance (infinity), where $U = 0$, with zero speed,
so $K = 0$. Its total energy there is zero. Energy is conserved, so the total energy at the surface
must also be zero:
\[ \frac{1}{2}mv_\mathrm{esc}^2 - \frac{GMm}{r} = 0 \quad\Rightarrow\quad v_\mathrm{esc}^2 = \frac{2GM}{r}
   \quad\Rightarrow\quad v_\mathrm{esc} = \sqrt{\frac{2GM}{r}} \]
\anspart{b}
\begin{crit}
  \critrow{Correctly calculates the escape velocity}{1}
\end{crit}
\sample
$v_\mathrm{esc} = \sqrt{\dfrac{2(\num{6.67e-11})(\num{6.42e23})}{\num{3.39e6}}} = \qty{5.0e3}{m.s^{-1}}$.

% ---- Q19
\ansq
\begin{crit}
  \critrow{States that the total energy is conserved, explains that $U$ is lowest (most negative)
    when closest, AND concludes that $K$, and so the speed, is greatest there}{3}
  \critrow{Gives two of these points}{2}
  \critrow{Gives one relevant point}{1}
\end{crit}
\sample
Only gravity acts on the comet, so its total energy $K + U$ stays the same. Its potential energy is
$U = -\dfrac{GMm}{r}$. When the comet is closest to the Sun, $r$ is smallest, so $U$ is the most
negative (lowest). For the total to stay the same, $K$ must be greatest there. So the comet's speed is
greatest when it is closest to the Sun. (This agrees with Kepler's second law.)

% ---- Q20
\ansq
\begin{crit}
  \critrow{Makes a clear judgement, supported by a correct calculation of $g$ at the ISS AND a
    correct explanation of apparent weightlessness (free fall, gravity provides the centripetal
    force, no normal force)}{5}
  \critrow{Gives a correct calculation AND a correct explanation, but the judgement is unclear or
    there is a minor error}{4}
  \critrow{Gives a correct calculation and some explanation, OR a good explanation without a
    calculation}{3}
  \critrow{Gives a correct calculation OR some relevant explanation}{2}
  \critrow{Gives one relevant point}{1}
\end{crit}
\sample
At the ISS, $r = \num{6.371e6} + \num{4.2e5} = \qty{6.791e6}{m}$, so
\[ g = \frac{GM}{r^2} = \frac{(\num{6.67e-11})(\num{6.0e24})}{(\num{6.791e6})^2} = \qty{8.7}{N.kg^{-1}} \]
This is about 89\% of the value at Earth's surface, so there is plenty of gravity at the ISS.

The astronauts float because they and the ISS are in free fall together. Gravity provides the
centripetal force that keeps them in orbit, and they all accelerate towards Earth at the same rate
($g$). The floor does not need to push on the astronauts, so there is no normal force and they feel
weightless (apparent weightlessness). Their horizontal speed (about \qty{7.7}{km.s^{-1}}) means they
keep falling around Earth instead of into it.

Judgement: the student's observation (the astronauts float) is correct, but the explanation is
wrong. Gravity is not zero in space; the floating is caused by free fall.

\end{document}
